% Generated draft: requires mathematical and source review.
\documentclass[11pt]{article}
\usepackage[margin=25mm]{geometry}
\usepackage{fontspec}
\setmainfont{Latin Modern Roman}
\usepackage{amsmath,amssymb,mathtools,unicode-math}
\setmathfont{Latin Modern Math}
\usepackage{enumitem,needspace,xurl,hyperref}
\setlength{\parindent}{0pt}
\setlength{\parskip}{5pt}
\setlength{\emergencystretch}{3em}
\newcommand{\N}{\mathbb{N}}
\newcommand{\Z}{\mathbb{Z}}
\newcommand{\Q}{\mathbb{Q}}
\newcommand{\R}{\mathbb{R}}
\newcommand{\C}{\mathbb{C}}
\newcommand{\F}{\mathbb{F}}
\newcommand{\A}{\mathbb{A}}
\newcommand{\PP}{\mathbb{P}}
\newcommand{\E}{\mathbb{E}}
\newcommand{\eps}{\varepsilon}
\newcommand{\dd}{\,\mathrm{d}}
\begin{document}
\section*{UnsolvedMath ID: 800803}
\textbf{Source problem code:} catalogue entry
\subsection{Definitions and mathematical statement}
A Lie groupoid \(\mathcal G\rightrightarrows M\) consists of a space \(M\) of objects and a space \(\mathcal G\) of arrows, with smooth source and target maps \(s,t:\mathcal G\to M\), units, inversion, and associative composition of composable arrows, satisfying the groupoid axioms. Its Lie algebroid has fibers and anchor
\[
A_x=\ker(\dd s)_{1_x},\qquad \rho_x=(\dd t)_{1_x}|_{A_x},\qquad x\in M,
\]
with the bracket induced by right-invariant vector fields tangent to the source fibers. The smooth category, including any infinite-dimensional analytic framework, must be specified in an elasticity application. A source-wise right-invariant metric is a metric on the tangent spaces of the fibers \(s^{-1}(x)\) for which groupoid right translations are isometries.

\textbf{Research problem (Khesin, Misiolek, Modin, 2022, Problem 2).} Extend the geodesic and Hamiltonian descriptions of Euler--Arnold equations on Lie groupoids to problems of elasticity theory. The intended task is to identify suitable elasticity configuration groupoids and associated Lie algebroids, together with geometric and Hamiltonian structures that describe the corresponding elasticity evolution equations. Where a geodesic description applies, the correspondence between the elasticity equations and the relevant geodesic equations must be established; a Hamiltonian description likewise requires a specified phase space, Poisson structure, Hamiltonian, and correspondence with the elasticity equations.

The supplied motivation reports that configuration spaces for multiphase fluids, ``homogenized'' vortex sheets, and generalized flows carry natural groupoid structures. It attributes to [IK18, IK22] the result that solutions of the Euler equations discussed in the source are precisely geodesics of an \(L^2\)-type source-wise right-invariant metric on the corresponding Lie groupoids of multiphase volume-preserving diffeomorphisms. The proposed extension replaces this hydrodynamical setting with elasticity applications, for which the source points to [mar-hugh, Kupf, JimLeon]. These are supplied background claims and reference labels, not independently verified results here.

\textbf{Unresolved scope.} The excerpt does not specify an elasticity model, material or spatial domain, constitutive law, boundary conditions, regularity class, configuration groupoid, or quantifiers over admissible elasticity problems. It also does not determine whether elastic potential energy should enter through a Hamiltonian or through an additional geometric construction. Consequently, this is a research-program statement rather than a fully specified existence or equivalence theorem; no universal claim over elasticity models is imposed.
\subsection{Short English statement}
Develop Lie groupoid formulations of elasticity equations that extend the geodesic and Hamiltonian formulations used in geometric hydrodynamics. The supplied record leaves the elasticity models and analytic hypotheses unspecified, so a precise theorem requires further choices and human mathematical review. Its classification as open is supplied registry metadata.
\subsection{Sources}
\begin{itemize}
\item Local UnsolvedMath registry, permanent ID \(800803\), source problem code \texttt{catalogue entry}, title: \emph{Geometric Hydrodynamics in Open Problems --- Problem 2}. The supplied record attributes the problem to Khesin, Misiolek, and Modin, with proposed year 2022, and labels it open; this status is not independently verified here.
\item Khesin, Misiolek, Modin, \emph{Geometric Hydrodynamics in Open Problems} (2022), Problem 2. Supplied source URL: \url{https://arxiv.org/abs/2205.01143}. The source record reports access on 2026-10-16 and extraction by \texttt{arxiv\_latex}; these are supplied provenance statements, not a claim of access or verification for this draft.
\item [IK18] and [IK22]: reference labels supplied for the hydrodynamical groupoid, Lie algebroid, geodesic, and Hamiltonian background. Full bibliographic entries are absent from the supplied record.
\item [mar-hugh], [Kupf], and [JimLeon]: reference labels supplied for applications of Lie groupoids to elasticity theory. Full bibliographic entries are absent from the supplied record.
\end{itemize}
\end{document}
